\section{Limitations}
\label{sec:limitations}

We explicitly acknowledge the following limitations:

\begin{enumerate}
    \item \textbf{Local approximation validity}: Our theory relies on $\eps_{\mathrm{sym}} < \delta$ and finite-step assumptions. Global geodesic interpolation and nonlinear large-deformation extrapolation are left as future work.

    \item \textbf{Pairwise parameter correspondence}: When comparing parameter vectors across different training runs (e.g., MAML vs independent training), neuron-wise comparison requires fixing permutation ambiguity. We use activation matching to establish correspondence; this preprocessing step adds computational overhead. Approximate acceleration strategies are provided in the appendix.

    \item \textbf{NTK diagnostic boundary}: SIREN with small hidden width (e.g., 32) deviates from infinite-width NTK assumptions. Our NTK analysis therefore serves only as a local curvature probe, not as an exact descriptor of network behavior. Finite-width low-rank curvature verification is provided in the appendix.

    \item \textbf{LIIF latent code separation}: The latent vector $z$ in LIIF's encoder-decoder architecture belongs to the pixel coordinate embedding space, not to the INR parameter orbit. We use LIIF as an architectural testbed to demonstrate generalization of the tangent initialization idea; the theoretical analysis of $\theta^*$ orbits does not depend on $z$.

    \item \textbf{Tangent initialization scope}: Our method is effective within the stable $\delta$ regime. Far from the training manifold (e.g., large extrapolation in scale or rotation), second-order curvature corrections or geodesic tracking should be combined.
\end{enumerate}

\input{figures/fig_limitations_schematic}
